\scanpage{6}
réduites, c'est un isom, cqfd.

\uline{Remarque} J'ignore si le corollaire reste valable sans supposer $A$ projectif sur $S$, tt comme j'ignore si $\underline{\mathrm{Pic}}_{A/S}$ existe dans ce cas, et s'il est \unsure{encore|ou non|même} simple sur $S$ le long de la section unité\ldots.

\uline{Corollaire 2} Si $\underline L$ est ample rel. \unsure{à|:} $S$, alors (et alors \unsure{seult|seult.}) $\varphi_{\underline L}:A\unsure{\text{---}|\longrightarrow} A'$ est un morphisme surjectif, \unsure{et même|étant|et} \inserted{\unsure{fini}} fidèlement plat ; son noyau est un groupe fini et plat sur $\unsure{S|s}$.
% slanted
\margin{énoncer sous forme de condition \unsure{néc.|nèc.} et suffisante sur $\underline L$.}

\uline{Corollaire} 3 $\underline{\mathrm{Pic}}^+_{\deleted{\unsure{X}}A/S}$ est ouvert \uline{et fermé}

[le ouvert est toujours le cas ; le fermé est spécial aux schémas abéliens].

En effet, si $\varphi_{\underline L}$ induit un hom surjectif en $s$, il induit un hom surjectif aux pts de $\overline s$.

\uline{Corollaire} 4 Existence de $\underline{\text{Nér}}^{\unsure{\varepsilon|\xi|S|\text{Sév}}}_{\deleted{\unsure{X}}A/S}$
% archive sheet number: 5

